\chapter{Considerações Parciais}

O presente trabalho propõe investigar e desenvolver uma solução para balanceamento de carga em Sistemas de Distribuição de Conteúdo, um componente crítico para garantir Qualidade de Experiência (QoE) ao usuário final. O trabalho identificou que estratégias convencionais de balanceamento, como \textit{Round-Robin} ou Seleção Aleatória, frequentemente ignoram o estado real dos recursos dos servidores, especialmente a memória, que é um fator de gargalo em sistemas cuja operação principal é a recuperação e entrega de dados oriundos de arquivos persistidos em memória secundária.

Diante dessa lacuna, foi elaborado um algoritmo de balanceamento probabilístico  que integra o monitoramento sistemático de memória principal e secundária como fator decisório. A proposta adapta o conceito de \textit{Load Interpretation}  para ajustar a distribuição de carga com base não apenas nas métricas de consumo de recursos (isto é, nesse caso, memória principal e secundária), mas também na idade dessas informações, evitando assim os efeitos adversos de dados desatualizados. A carga de cada servidor ($L_i$) é calculada como uma média ponderada do consumo normalizado de memória RAM ($LM_i$) e da carga de disco ($LD_i$), conforme a Equação \ref{eq:limemory}. A probabilidade de um servidor ser escolhido ($P_i$) é então determinada dinamicamente, de modo a equalizar a carga projetada entre todos os nós do sistema.

A arquitetura do sistema para a fase experimental foi definida para operar em um ambiente controlado e reproduzível com contêineres Docker , utilizando servidores MPEG-DASH, um \textit{broker} MQTT para a comunicação assíncrona dos dados de monitoramento  e um balanceador de carga em Rust. A validação da proposta será realizada por meio de uma análise comparativa rigorosa com os algoritmos \textit{Round-Robin} e Seleção Aleatória, sob cenários de alta carga, utilizando métricas de QoE como o número e a duração de travamentos do vídeo.

Com a fundamentação teórica e a arquitetura propostas, o trabalho está preparado para a fase de implementação e validação empírica. As atividades subsequentes estão organizadas para construir, testar e avaliar sistematicamente cada componente, culminando na análise final que determinará a eficácia do algoritmo proposto em otimizar a distribuição de conteúdo de maneira inteligente e resiliente.

\section{Próximos passos}

Esta seção descreve as atividades que serão feitas a seguir, bem como o planejamento de execução.

\begin{enumerate}
    \item Configuração e testes do sistema gerenciador de filas (\textit{broker} MQTT);
    \item Implementação e configuração inicial do servidor MPEG-DASH
    \begin{enumerate}
        \item[2.1.] Configuração e testes com o \textit{player} enviando diretamente requisições para o servidor;
        \item[2.2.] Instrumentação do servidor, enviando publicações periódicas ao serviço gerenciador de filas com as informações necessárias;
    \end{enumerate}
    \item Implementação do balanceador de carga;
    \begin{enumerate}
        \item[3.1.] Criação da arquitetura de código flexível para escolher o algoritmo usado no balanceamento;
        \item[3.2.] Implementação e teste do balanceador com servidores falsos e algoritmos simples;
        \item[3.3.] Implementação do algoritmo proposto no balanceador de carga;
        \item[3.4.] Testes do balanceador de carga ligado nos servidores reais;
    \end{enumerate}
    \item Criação da aplicação cliente geradora de sobrecarga no sistema;
    \begin{enumerate}
        \item[4.1.] Teste da aplicação com um servidor MPEG-DASH público;
        \item[4.2.] Teste da aplicação ligada em um dos servidores criados (diretamente - sem balanceador de carga);
        \item[4.3.] Teste da aplicação ligada no balanceador de carga com todos os servidores configurados;
    \end{enumerate}
    \item Configuração do \textit{player} simples no navegador;
    \item Execução dos testes de qualidade com o \textit{player} simples;
    \item Elaboração dos gráficos e descrição dos resultados finais;
\end{enumerate}

\section{Cronograma}

\vspace{0.5cm}
\begin{center} % Centraliza a tabela
    \begin{tabular}{|c|c|c|c|c|c|c|c|}
    \hline
    \multirow{2}{*}{\textbf{\small{Atividades}}} &
    \multicolumn{7}{|c|}{\textbf{\small{Meses}}} \\
    \cline{2-8} % Ajustado para 7 meses + 1 coluna de atividades
    &\textbf{Jun} & \textbf{Jul} & \textbf{Ago} & \textbf{Set} &
    \textbf{Out} & \textbf{Nov} & \textbf{Dez} \\
    \hline
    \textbf{\small{1}} & \cellcolor{gray} & \cellcolor{gray} & & & & & \\
    \hline
    \textbf{\small{2}} & \cellcolor{gray} & \cellcolor{gray} & & & & & \\
    \hline
    \textbf{\small{3}} & & & \cellcolor{gray} & \cellcolor{gray} & & & \\
    \hline
    \textbf{\small{4}} & & & & & \cellcolor{gray} & & \\
    \hline
    \textbf{\small{5}} & & & & & \cellcolor{gray} & & \\
    \hline
    \textbf{\small{6}} & & & & & & \cellcolor{gray} & \cellcolor{gray} \\
    \hline
    \textbf{\small{7}} & & & & & & \cellcolor{gray} & \cellcolor{gray} \\
    \hline
    \end{tabular}
\end{center}

\vspace{0.5cm}
